\documentclass[12pt]{article}

\usepackage[utf8]{inputenc}
\usepackage[T1]{fontenc}
\usepackage{lmodern}
\usepackage[a4paper,margin=1in]{geometry}
\usepackage{amsmath,amssymb,amsthm}

\begin{document}

\section*{Statement}

Each of two equal-area squares is partitioned into 100 equal-area parts. Prove that the two squares can be placed one over the other and pierced in 100 points so that every part of each square is pierced.

\section*{Solution}

Let the parts of the first square be \(A_1, A_2, \dots, A_{100}\), and the parts of the second square be \(B_1, B_2, \dots, B_{100}\). Let \(S\) denote the area of each square. Since each square is partitioned into 100 equal-area parts,
\[
S(A_i)=S(B_j)=\frac{S}{100}
\qquad
\text{for all } i,j.
\]

We construct a bipartite graph whose left part consists of the vertices
\[
A_1,\dots,A_{100},
\]
and whose right part consists of the vertices
\[
B_1,\dots,B_{100}.
\]
We connect \(A_i\) and \(B_j\) by an edge whenever
\[
S(A_i\cap B_j)>0.
\]

We verify Hall's condition. Take any subset \(I\subseteq \{1,\dots,100\}\), and denote
\[
U_I=\bigcup_{i\in I} A_i.
\]
Let \(J(I)\) be the set of all indices \(j\in \{1,\dots,100\}\) such that
\[
S(U_I\cap B_j)>0.
\]
Then
\[
U_I\subseteq \bigcup_{j\in J(I)} B_j.
\]
Therefore,
\[
S(U_I)\le \sum_{j\in J(I)} S(B_j)=|J(I)|\cdot \frac{S}{100}.
\]
On the other hand, the sets \(A_i\) are pairwise disjoint and each has area \(S/100\), so
\[
S(U_I)=|I|\cdot \frac{S}{100}.
\]
Hence
\[
|I|\cdot \frac{S}{100}\le |J(I)|\cdot \frac{S}{100},
\]
and therefore
\[
|I|\le |J(I)|.
\]
Thus Hall's condition holds:
\[
\forall I\subseteq \{1,\dots,100\}\qquad |J(I)|\ge |I|.
\]

By Hall's theorem, there exists a perfect matching in this bipartite graph. Equivalently, there exists a permutation \(\pi\) of \(\{1,\dots,100\}\) such that
\[
S(A_i\cap B_{\pi(i)})>0
\qquad
\text{for every } i=1,\dots,100.
\]

For each \(i\), choose a point
\[
x_i\in A_i\cap B_{\pi(i)}.
\]
Then each point \(x_i\) pierces the part \(A_i\) of the first square and the part \(B_{\pi(i)}\) of the second square. Since \(\pi\) is a permutation, every part \(A_i\) and every part \(B_j\) is pierced by one of the points \(x_1,\dots,x_{100}\).

Therefore, 100 points can be chosen so that every part of each square is pierced.

\end{document}
